\section{Conclusion}


This paper presented the Unified Dual-Mask Physical Model, a comprehensive framework designed to address the limitations of Epipolar Plane Image (EPI) assumptions in light field depth estimation, particularly for non-Lambertian surfaces and complex textures. By integrating physical reflectance properties with geometric feature analysis, the proposed approach systematically mitigates the slope ambiguity and structural degradation typically encountered in mixed and urban environments. 

The specific contributions and findings are summarized as follows:
1. \textbf{Three-Layer Angular Signal Decomposition and Geometric Material Perception}: We developed a material perception mechanism that extracts geometric features to replace traditional frequency-domain representations. This design overcomes the resolution limits of discrete angular sampling and provides a robust routing basis, effectively identifying the physical causes of EPI assumption failures on non-Lambertian surfaces.
2. \textbf{Dual-Mask Routing Mechanism for EPI Structure Preservation}: We introduced a dual-mask architecture that dynamically routes and processes EPI features. This mechanism reduces slope blurring caused by complex textures and structural disruptions induced by specular reflections, achieving an overall Mean Absolute Error (MAE) of 0.133 and a Mixed (Urban) scene MAE of 0.081.
3. \textbf{Domain-Balanced Sampling and Physical Boundary Validation}: We formulated a domain-balanced training strategy that maximizes generalization under severe data scarcity, specifically addressing the constraint of only four non-Lambertian training scenes. Through 132 experiments across 16 research directions, we established the terminal physical boundaries of EPI-based unified models, proving that architectural modifications cannot overcome the dual barriers of physical assumption breakdown and extreme data deficiency in non-Lambertian and complex-texture Lambertian scenarios.

The established physical boundaries and the proposed dual-mask routing paradigm provide a definitive theoretical and empirical reference for computational photography. These findings indicate that advancing light field 3D reconstruction necessitates a paradigm shift from pure EPI reliance toward hybrid representations capable of modeling complex light transport.


The fundamental bottleneck in light field depth estimation lies not merely in algorithmic capacity, but in the intrinsic physical coupling of material reflectance and geometric depth. By decoupling these signals through a three-layer angular decomposition, the proposed framework establishes a new paradigm for solving ill-posed inverse problems in computational photography, shifting the focus from purely data-driven heuristics to physically grounded system design.

Unlike conventional frequency-domain approaches that implicitly assume dense angular sampling, our system-level design explicitly acknowledges the hardware constraints of commercial plenoptic cameras. Under low angular resolutions (e.g., $5 \times 5$ or $9 \times 9$ discrete sampling), 2D-DFT suffers from severe spectral aliasing due to the violation of the Nyquist-Shannon sampling theorem. By shifting material perception to the spatial domain via geometric angular features, we bypass this hardware-imposed resolution bottleneck. The exceptionally high Cohen's $d$ effect size ($0.983$) for the angular gradient feature empirically validates that spatial-domain geometry captures reflectance discrepancies more faithfully than aliased frequency spectra, ensuring robust performance without requiring computationally prohibitive super-resolution preprocessing.

Furthermore, the integration of a Random Forest (RF) classifier to generate the dual-mask routing mechanism introduces a crucial architectural advantage: the strict decoupling of material classification from depth regression. End-to-end joint optimization of material and depth often leads to gradient conflicts and ambiguous soft masks, particularly in texture-less or highly specular urban scenes. By employing an RF classifier, we achieve a computationally lightweight, hard-boundary physical mask. This explicit, interpretable routing acts as a deterministic prior that stabilizes the subsequent deep feature extraction pipeline, preventing the network from overfitting to spurious Epipolar Plane Image (EPI) slope correlations while adding negligible inference latency to the overall system.

Despite these systemic advancements, the current physical model operates under the assumption of local illumination consistency within the three-layer decomposition. In scenarios dominated by complex global illumination effects, such as inter-reflections or caustics, the residual term may conflate secondary light bounces with intrinsic material properties, potentially degrading the mask accuracy. Additionally, while the discrete EPI-based routing is highly efficient for multi-view stereo pipelines, it inherently limits the continuous volumetric representation of highly scattering or translucent materials.

Future work will explore integrating this physical dual-mask prior into continuous 3D representations, such as Neural Radiance Fields (NeRF) or 3D Gaussian Splatting. By embedding the geometric angular features directly into the volumetric rendering equation:
\begin{equation}
\label{eq:eq4}
L_o(\mathbf{x}, \omega_o) = \int_{\Omega} f_r(\mathbf{x}, \omega_i, \omega_o) L_i(\mathbf{x}, \omega_i) (\omega_i \cdot \mathbf{n}) d\omega_i
\end{equation}
we anticipate achieving unified, view-consistent depth and material estimation even under extreme non-Lambertian conditions. This extension would ultimately bridge the gap between discrete light field processing and continuous physical scene reconstruction, enabling real-time, physically plausible 3D perception for autonomous navigation systems.

Despite the demonstrated efficacy of the Unified Dual-Mask Physical Model in decoupling material reflectance from geometric depth, several inherent limitations remain that outline promising directions for future research. 

First, the current material perception mechanism relies on a Random Forest classifier to generate hard binary masks for Lambertian and non-Lambertian regions. While this design provides robust and computationally efficient routing priors, hard boundaries can inadvertently introduce depth discontinuities or structural artifacts at material transition zones, such as the penumbra of specular highlights or translucent-to-opaque gradients. Future work will investigate differentiable soft routing mechanisms and uncertainty-aware mask generation to ensure smooth, physically plausible depth transitions across mixed-material boundaries without sacrificing inference speed.

Second, the proposed three-layer angular signal decomposition fundamentally assumes a static scene with relatively stable ambient illumination represented by the DC component. In highly dynamic environments—such as fluctuating water surfaces, moving specular reflections, or extreme low-light and overexposed conditions—the linear decoupling assumption may degrade. Extending the physical model to incorporate temporal consistency through video light fields could provide crucial supplementary constraints to disambiguate dynamic non-Lambertian effects from true geometric parallax.

Third, although geometric angular features successfully bypass the spectral aliasing of 2D-DFT at standard low angular resolutions (e.g., $5 \times 5$ or $9 \times 9$), the computation of angular gradients and correlations becomes inherently ill-posed at ultra-sparse angular samplings (e.g., $2 \times 2$ or $3 \times 3$). At such extreme sparsity, the discrete approximation of the angular derivative loses its physical meaning. To address this, future iterations could integrate learned angular super-resolution or implicit neural representations to synthesize dense, continuous angular features prior to mask generation, thereby extending the operational envelope of the physical model to highly constrained hardware setups.

Finally, while this framework significantly advances 2.5D depth estimation, integrating the dual-mask physical priors into modern volumetric rendering paradigms, such as Neural Radiance Fields (NeRF) or 3D Gaussian Splatting, represents a critical next step. Current neural rendering methods often struggle with view-dependent non-Lambertian effects, leading to severe floaters and specular flickering. By embedding our material-specific geometric routing directly into the radiance field optimization pipeline, we can explicitly regularize the view-dependent color networks. This integration would not only resolve rendering artifacts but also bridge the fundamental gap between robust geometric depth estimation and photorealistic, physically grounded 3D scene reconstruction.